\documentclass[12pt,a4paper]{article}

\usepackage[utf8]{inputenc}
\usepackage{amsmath,amssymb,amsthm}
\usepackage{geometry}
\usepackage{setspace}
\usepackage{natbib}
\usepackage{hyperref}
\usepackage{booktabs}

\geometry{margin=1in}
\onehalfspacing

\newtheorem{proposition}{Proposition}
\newtheorem{corollary}{Corollary}
\newtheorem{remark}{Remark}
\newtheorem{definition}{Definition}

\title{\textbf{Debt Denomination and Fiscal Valuation:\\
Nominal, Floating-Rate, and Inflation-Indexed Government Liabilities}}

\author{}
\date{\today}

\begin{document}

\maketitle

\begin{abstract}
This note develops a fiscal valuation model with three classes of government liabilities:
nominal fixed-rate debt, floating-rate debt, and inflation-indexed debt. Standard
formulations of the Fiscal Theory of the Price Level emphasize nominal government
liabilities, whose real value can be reduced by inflation. Many governments, however,
issue large quantities of inflation-indexed and floating-rate debt. These instruments
alter the margins through which fiscal valuation can be restored.

We derive a generalized government valuation equation with multiperiod nominal,
floating-rate, and indexed bonds. We then define the valuation elasticity of each debt
instrument with respect to the price level. Nominal fixed-rate debt has the highest
valuation elasticity because surprise inflation directly erodes its real value.
Inflation-indexed debt has zero direct valuation elasticity because its payoff is defined
in real terms. Floating-rate debt occupies an intermediate position: although its real
principal is eroded by inflation, monetary policy feedback raises future debt-service
costs. The model therefore delivers a three-way ranking of fiscal inflation incentives:
nominal fixed-rate debt generates the strongest price-level valuation channel,
floating-rate debt an intermediate one, and indexed debt the weakest one.

The note provides a theoretical foundation for empirical work using debt denomination
to distinguish fiscal valuation mechanisms from resource-constraint explanations of
inflation.
\end{abstract}

\newpage

\section{Introduction}

The Fiscal Theory of the Price Level (FTPL) views the price level as part of the
valuation mechanism for government liabilities. In its simplest form, the real value of
nominal government debt must equal the expected present value of future primary
surpluses. If fiscal backing is insufficient, the price level can adjust so that the real
value of nominal liabilities falls.

This standard intuition is sharpest when all public debt is nominal and fixed in currency
units. Modern government debt portfolios, however, are more complex. In several
countries, and especially in Brazil, the Treasury issues not only nominal fixed-rate bonds,
but also inflation-indexed bonds and floating-rate securities. These instruments differ in
the extent to which inflation can reduce their real fiscal burden.

This note develops a compact theoretical framework that incorporates these three debt
classes into the FTPL valuation equation. The purpose is not to provide a full quantitative
model of debt management, but to isolate a simple implication: debt denomination
determines the strength of the price-level valuation channel.

The key result is a ranking of valuation elasticities. Nominal fixed-rate debt has the
highest valuation elasticity with respect to the price level. Inflation-indexed debt has no
direct valuation elasticity. Floating-rate debt lies between the two because inflation erodes
the real value of outstanding nominal principal, but higher short-term interest rates feed
back into future debt-service costs.

The contribution is therefore theoretical and clarifying. It shows how the FTPL
valuation logic extends beyond nominal debt and provides a foundation for empirical
tests based on the denomination composition of government liabilities.

\section{Household Environment and Asset Pricing}

Consider a discrete-time endowment economy populated by an infinitely lived
representative household. Preferences are given by
\begin{equation}
\mathbb{E}_0
\sum_{t=0}^{\infty}
\beta^t
U\left(C_t,\frac{M_t}{P_t}\right),
\qquad
\beta \in (0,1),
\end{equation}
where $C_t$ is consumption, $M_t/P_t$ denotes real money balances, and $P_t$ is the
aggregate price level.

Let
\begin{equation}
m_{t,t+1}
\equiv
\beta
\frac{
U_C(C_{t+1},M_{t+1}/P_{t+1})
}{
U_C(C_t,M_t/P_t)
}
\end{equation}
denote the one-period real stochastic discount factor. For $\tau \geq 1$, define the
multiperiod real stochastic discount factor
\begin{equation}
m_{t,t+\tau}
\equiv
\prod_{s=0}^{\tau-1} m_{t+s,t+s+1}.
\end{equation}

The government issues three types of liabilities.

\begin{enumerate}
\item \textbf{Nominal fixed-rate bonds.} A nominal bond issued at time $t$ promises
a fixed amount of currency in the future. Its real payoff is therefore exposed to
future inflation.


\item \textbf{Inflation-indexed bonds.} An indexed bond promises payments defined
in units of the consumption good. Its real payoff is protected from inflation by
construction.

\item \textbf{Floating-rate bonds.} A floating-rate bond promises nominal payments
indexed to the short-term policy rate. Its payoff is nominal, but future coupons
adjust with monetary policy.


\end{enumerate}

For a one-period nominal fixed-rate bond with gross nominal return $1+i_t^N$, the Euler
equation is
\begin{equation}
1=\mathbb{E}_t
\left[
m_{t,t+1}
(1+i_t^N)
\frac{P_t}{P_{t+1}}
\right].
\label{eq:euler_nominal}
\end{equation}

For a one-period inflation-indexed bond with gross real return $1+i_t^R$, the Euler
equation is
\begin{equation}
1=\mathbb{E}_t
\left[
m_{t,t+1}
(1+i_t^R)
\right].
\label{eq:euler_real}
\end{equation}

For a one-period floating-rate bond whose payoff is indexed to the within-period policy
rate $i_t$, the Euler equation is
\begin{equation}
1=\mathbb{E}_t
\left[
m_{t,t+1}
(1+i_t)
\frac{P_t}{P_{t+1}}
\right].
\label{eq:euler_float}
\end{equation}

The difference between \eqref{eq:euler_nominal} and \eqref{eq:euler_real} is the
standard Fisher distinction between nominal and real payoffs. The floating-rate
condition \eqref{eq:euler_float} is closer to the nominal bond equation because the
payoff is denominated in currency. However, unlike fixed-rate nominal bonds, the
floating-rate payoff adjusts with the policy rate. This feature is central for fiscal
valuation.

\section{Government Valuation with Three Debt Classes}

Let the government issue portfolios of nominal fixed-rate, floating-rate, and
inflation-indexed bonds at multiple maturities. Let $B_t^{N,(\tau)}$ denote the nominal
face value of fixed-rate debt maturing in $\tau$ periods, $B_t^{F,(\tau)}$ the nominal
value of floating-rate debt, and $B_t^{R,(\tau)}$ the real face value of indexed debt.

Let $Q_t^{N,(\tau)}$, $Q_t^{F,(\tau)}$, and $Q_t^{R,(\tau)}$ denote the corresponding
time-$t$ prices. The real market value of outstanding government liabilities is

\begin{equation}
\mathcal{V}_t
=
\underbrace{
\frac{1}{P_t}
\sum_{\tau\geq 1}
Q_t^{N,(\tau)}
B_t^{N,(\tau)}
}_{\mathcal{V}_t^N}
+
\underbrace{
\frac{1}{P_t}
\sum_{\tau\geq 1}
Q_t^{F,(\tau)}
B_t^{F,(\tau)}
}_{\mathcal{V}_t^F}
+
\underbrace{
\sum_{\tau\geq 1}
Q_t^{R,(\tau)}
B_t^{R,(\tau)}
}_{\mathcal{V}_t^R}.
\label{eq:market_value}
\end{equation}


Nominal fixed-rate and floating-rate debt enter in real terms after division by $P_t$.
Inflation-indexed debt enters directly in real terms because its payoff is already
defined in consumption units.

Let $S_t$ denote the real primary surplus. Ignoring seigniorage for notational
simplicity, the intertemporal government valuation equation is

\begin{equation}
\boxed{
\mathcal{V}_t^N+\mathcal{V}_t^F+\mathcal{V}_t^R
=
\mathbb{E}_t
\sum_{\tau=0}^{\infty}
m_{t,t+\tau}
S_{t+\tau}
\equiv
PV_t(S)
}
\label{eq:valuation_general}
\end{equation}


Equation \eqref{eq:valuation_general} is the natural extension of the standard FTPL
valuation equation to a debt portfolio with nominal fixed-rate, floating-rate, and
inflation-indexed liabilities.


\begin{proposition}[Fiscal valuation with heterogeneous debt]
Suppose the government's transversality condition holds:
\begin{equation}
\lim_{T\to\infty}
\mathbb{E}_t
\left[
m_{t,T}\mathcal{V}_T
\right]
=
0.
\end{equation}
Then the real market value of government liabilities satisfies the valuation equation
\eqref{eq:valuation_general}.
\end{proposition}


\begin{proof}
Starting from the real government flow budget constraint, the real value of liabilities
carried into the next period equals current primary deficits plus the real value of
maturing obligations rolled into new liabilities. Iterating the budget constraint forward
and imposing the transversality condition eliminates the terminal value of debt. The
remaining object is the expected discounted value of future primary surpluses. Because
the current liability portfolio is marked to market, the left-hand side is the sum of the
real market values of nominal fixed-rate, floating-rate, and inflation-indexed liabilities.
\end{proof}

\section{Valuation Elasticities}

The three debt classes differ in the extent to which a change in the price level reduces
their real fiscal burden.


\begin{definition}[Valuation elasticity]
For debt class $j \in \{N,F,R\}$, define its price-level valuation elasticity as
\begin{equation}
\kappa_t^j
\equiv
-
\frac{\partial \log \mathcal{V}_t^j}{\partial \log P_t},
\label{eq:kappa_def}
\end{equation}
where the derivative is evaluated holding fixed the contractual quantity of debt and
allowing for the relevant pricing response of that instrument.
\end{definition}


For nominal fixed-rate debt, the payoff is predetermined in currency units. Holding
nominal bond prices fixed on impact, a rise in the price level reduces the real value of
the claim one-for-one:
\begin{equation}
\kappa_t^N = 1.
\label{eq:kappa_nominal}
\end{equation}

For inflation-indexed debt, the payoff is defined in units of the consumption good. Its
real value is not directly reduced by a change in the price level:
\begin{equation}
\kappa_t^R = 0.
\label{eq:kappa_real}
\end{equation}

Floating-rate debt is intermediate. Its principal is nominal, so inflation has a direct
debt-eroding effect. However, because coupons reset with the short-term policy rate,
inflation also raises future debt-service costs when monetary policy reacts to inflation.
Let the effective valuation elasticity of floating-rate debt be written as

\begin{equation}
\kappa_t^F
=
1-\chi_t,
\label{eq:kappa_float}
\end{equation}

where $\chi_t$ denotes the fraction of direct inflation relief offset by higher expected
future floating-rate payments. Under an active monetary regime in which the policy rate
responds positively to inflation, $\chi_t>0$. If the feedback does not more than offset
the direct inflation relief, $\chi_t<1$.

Thus,
\begin{equation}
0 < \kappa_t^F < 1.
\label{eq:kappa_float_bounds}
\end{equation}


\begin{proposition}[Ranking of valuation elasticities]
\label{prop:kappa_ranking}
Under an active monetary regime in which floating-rate coupons respond positively but
not explosively to inflation, the valuation elasticities satisfy
\begin{equation}
\kappa_t^N > \kappa_t^F > \kappa_t^R,
\qquad
\text{with}
\qquad
\kappa_t^N=1,
\quad
0<\kappa_t^F<1,
\quad
\kappa_t^R=0.
\label{eq:kappa_ranking}
\end{equation}
\end{proposition}



\begin{proof}
The nominal fixed-rate claim is fixed in currency units, so its real value falls
one-for-one with the price level, implying $\kappa_t^N=1$. The indexed claim is fixed in
real units, so its real value is invariant to the price level, implying $\kappa_t^R=0$.
The floating-rate claim is nominal, so it has the same direct price-level exposure as
nominal fixed-rate debt. However, its coupon payments adjust with the policy rate.
Under an active monetary regime, higher inflation raises expected future policy rates
and therefore increases future floating-rate debt service. This offsets part of the direct
inflation relief. If the offset is partial, then $0<\kappa_t^F<1$.
\end{proof}

\begin{remark}
The intermediate position of floating-rate debt is not a purely contractual property.
Unlike nominal fixed-rate debt and indexed debt, whose valuation elasticities are pinned
down directly by their payoff definitions, the floating-rate elasticity depends on the
monetary-policy regime. This is why floating-rate debt provides a sharper but more
regime-dependent implication of fiscal valuation.
\end{remark}

\section{Fiscal Inflation Incentives}

Consider an unbacked fiscal valuation shock associated with the issuance of debt type
$j$. Let the shock be normalized by market value. Suppose the component of the shock
that is not absorbed by revisions in expected future surpluses is transmitted through the
price-level valuation channel with strength proportional to $\kappa_t^j$.

A convenient reduced-form representation of the inflation response is

\begin{equation}
IRF_{\pi}^j(h)
=
\kappa_t^j
(1-\rho_t^j)
\Theta_h,
\label{eq:irf_kappa}
\end{equation}

where $\rho_t^j\in[0,1]$ denotes the residual fiscal-backing share of the shock and
$\Theta_h>0$ is a common propagation term.

Equation \eqref{eq:irf_kappa} separates two objects. The valuation elasticity
$\kappa_t^j$ is determined by debt denomination and the monetary regime. The residual
backing term $\rho_t^j$ captures how much of the shock is absorbed through changes in
expected future primary surpluses.

\begin{proposition}[Debt-denomination ranking]
Suppose residual fiscal backing is common across debt types:
\begin{equation}
\rho_t^N=\rho_t^F=\rho_t^R\equiv \rho_t<1.
\label{eq:common_rho}
\end{equation}
Then the inflation responses inherit the valuation-elasticity ranking:
\begin{equation}
|IRF_{\pi}^N(h)|
>
|IRF_{\pi}^F(h)|
>
|IRF_{\pi}^R(h)|.
\label{eq:irf_ranking}
\end{equation}
\end{proposition}


\begin{proof}
Under \eqref{eq:common_rho},
\begin{equation}
IRF_{\pi}^j(h)
=
\kappa_t^j(1-\rho_t)\Theta_h.
\end{equation}
Since $(1-\rho_t)\Theta_h>0$, the ranking of inflation responses is identical to the
ranking of $\kappa_t^j$. The result follows from Proposition~\ref{prop:kappa_ranking}.
\end{proof}



\begin{remark}
The ranking does not identify the level of fiscal backing. It identifies the relative
ordering of the valuation elasticities $\kappa_t^N$, $\kappa_t^F$, and $\kappa_t^R$:
the differential capacity of debt instruments to transmit fiscal valuation shocks through
the price level.
\end{remark}

\section{Robustness to Residual-Backing Heterogeneity}

Common residual backing is sufficient but not necessary for the ranking in
\eqref{eq:irf_ranking}. To characterize the weaker condition, define
\begin{equation}
\phi_t^j
\equiv
\kappa_t^j(1-\rho_t^j).
\end{equation}

The ranking survives whenever
\begin{equation}
\phi_t^N>\phi_t^F>\phi_t^R.
\label{eq:phi_ranking}
\end{equation}

Equivalently,

\begin{align}
\kappa_t^N\rho_t^N-\kappa_t^F\rho_t^F
&<
\kappa_t^N-\kappa_t^F,
\label{eq:region_nf}
\\
\kappa_t^F\rho_t^F-\kappa_t^R\rho_t^R
&<
\kappa_t^F-\kappa_t^R.
\label{eq:region_fr}
\end{align}


Equations \eqref{eq:region_nf} and \eqref{eq:region_fr} define an open convex region
of residual-backing heterogeneity under which the denomination ranking remains valid.

\begin{corollary}[Common residual backing is sufficient but not necessary]
The common-backing case $\rho_t^N=\rho_t^F=\rho_t^R=\rho_t<1$ lies inside the region
defined by \eqref{eq:region_nf}--\eqref{eq:region_fr}. Therefore, common residual
backing is sufficient but not necessary for the inflation-ranking restriction.
\end{corollary}

\begin{proof}
Substitute $\rho_t^N=\rho_t^F=\rho_t^R=\rho_t$ into
\eqref{eq:region_nf}. The condition becomes
\begin{equation}
\rho_t(\kappa_t^N-\kappa_t^F)
<
\kappa_t^N-\kappa_t^F,
\end{equation}
which holds whenever $\rho_t<1$. The same argument applied to
\eqref{eq:region_fr} gives
\begin{equation}
\rho_t(\kappa_t^F-\kappa_t^R)
<
\kappa_t^F-\kappa_t^R,
\end{equation}
which also holds whenever $\rho_t<1$.
\end{proof}

This robustness region clarifies the role of the assumptions. Violations of common
residual backing need not invalidate the ranking. They matter only if residual backing
differs across instruments strongly enough to overturn the ordering of
$\kappa_t^j(1-\rho_t^j)$.

A particularly robust comparison is between nominal fixed-rate and inflation-indexed
debt. Since $\kappa_t^N=1$ and $\kappa_t^R=0$, the nominal--indexed ranking is pinned
down by contract and does not depend on the monetary regime. The full three-way
nominal--floating--indexed ranking is sharper, but it is also more dependent on the
policy-rate feedback governing $\kappa_t^F$.

\section{Valuation Variables}

Indexed debt cannot be directly devalued through inflation. Therefore, an unbacked
increase in indexed liabilities must be absorbed through other valuation margins:
expected future primary surpluses, real discount rates, liquidity premia, or sovereign
risk premia.

This observation does not imply a strict theorem for any single asset-price variable.
For example, sovereign CDS spreads include default risk, currency risk, global risk
appetite, and liquidity components. Nevertheless, the model suggests a directional
prediction: shocks to indexed debt should place relatively greater pressure on valuation
variables than shocks to nominal fixed-rate debt.

Thus, if $r_t$ denotes a real yield and $s_t$ a sovereign spread, the model suggests
\begin{equation}
|IRF_r^R(h)|
\gtrsim
|IRF_r^F(h)|
\gtrsim
|IRF_r^N(h)|,
\end{equation}
and similarly for sovereign spreads, subject to the caveat that empirical spreads need
not map one-for-one into the theoretical valuation margin.

\section{Conclusion}

This note extends the fiscal valuation logic of the FTPL to a government debt portfolio
containing nominal fixed-rate, floating-rate, and inflation-indexed liabilities. The key
object is the valuation elasticity of each instrument with respect to the price level.

Nominal fixed-rate debt has the highest valuation elasticity because inflation directly
reduces its real value. Inflation-indexed debt has zero direct valuation elasticity because
its payoff is defined in real terms. Floating-rate debt lies between the two because the
direct erosion of real nominal principal is offset by higher future debt-service costs when
monetary policy responds to inflation.

The resulting ranking,
\begin{equation}
\kappa_t^N>\kappa_t^F>\kappa_t^R,
\end{equation}
provides a theoretical foundation for empirical tests using debt denomination. More
generally, it shows that the composition of government liabilities matters for the
mechanism through which fiscal imbalances are transmitted to prices and asset values.

\bibliographystyle{aea}
\bibliography{refs}

\end{document}
